\chapter{Exact operations, orthogonality, and stationary elimination}
\label{ch:exact-concepts}

The proof uses two kinds of elimination. An exact linear contraction
replaces a finite part of the tangle calculation by the same linear
map written more economically. Stationary elimination replaces
internal variables of an analytic phase by their critical values.
The first is an identity at every level; the second requires a
nondegenerate critical point and an estimate for the original sum.
This chapter records the elementary facts needed to keep these
operations and their estimates separate.

\section{The coefficient space and its actual pairing}

Let $V$ be a finite-dimensional complex inner-product space. We take
the inner product to be conjugate-linear in its first argument.
Choose vectors $v_1,\ldots,v_m$ spanning a subspace of $V$, and let
\[
 B:\CC^m\longrightarrow V,\qquad Bc=\sum_{j=1}^m c_jv_j.
\]
The spanning vectors need not be independent or orthogonal. The
adjoint is denoted by $B^*$; in orthonormal coordinates it is the
conjugate transpose. Define the Gram matrix $G=B^*B$.

\begin{proposition}[The transported coefficient pairing]
\label{prop:coefficient-gram}
For all $c,d\in\CC^m$,
\[
 \langle Bc,Bd\rangle=c^*Gd,\qquad
 \|Bc\|^2=c^*Gc,\qquad
 \ker G=\ker B.
\]
Consequently $G$ defines a positive definite pairing on
$\CC^m/\ker B$. If new coefficients are related by $c=Sd$, their
Gram matrix is $S^*GS$.
\end{proposition}
\begin{proof}
The adjoint identity gives
$\langle Bc,Bd\rangle=\langle c,B^*Bd\rangle$.
Thus $G$ is positive semidefinite. If $Gc=0$, its quadratic form
vanishes and $Bc=0$; the converse follows from $Gc=B^*(Bc)$.
This also shows that the form descends to the stated quotient and
has no remaining null vector. Substituting $c=Sd$ gives the last
formula.
\end{proof}

In particular, replacing $G$ by the identity is justified only for
orthonormal spanning vectors. Redundant coefficient directions have
zero energy because they represent the zero vector, not because a
later estimate is expected to cancel them.

A closing functional is additional data. For a fixed complex-linear
map $\ell:V\to\CC$,
\[
 |\ell(Bc)|\le \|\ell\|\,\sqrt{c^*Gc}.
\]
This is an upper bound. It gives no positive lower bound on
$|\ell(Bc)|$: a nonzero vector may lie in $\ker\ell$. For example,
the unit vector $(1,1)/\sqrt2$ is killed by
$\ell(z_1,z_2)=(z_1-z_2)/\sqrt2$.
The proof of nonvanishing must therefore concern the original
closing functional, not just the size of a coefficient vector.

In the Newton coordinates used for the tangles, the closing has the
bilinear form $a^{\mathsf T}Jb$, where $J$ reverses the coefficient
order. With the ordinary Hermitian coefficient norm,
$J^*J=I$, so
\[
 |a^{\mathsf T}Jb|\le\|a\|_2\|b\|_2.
\]
The equality with the physical cap and the normalization of these
coordinates are the exact finite identities cited in
Chapter~\ref{ch:conway-current}. The bilinear closing itself is not
a Hermitian inner product.

\section{Contracting a linear part of the diagram}

Suppose a sequence of already specified pieces acts by linear maps
\[
 V_0\xrightarrow{T_1}V_1\xrightarrow{T_2}\cdots
       \xrightarrow{T_s}V_s.
\]
Replacing these pieces by $T=T_s\cdots T_1$ is exact. We call such
a sequence a \emph{linear corridor}. A product operation with one
fixed input is linear in its other input and may occur in a
corridor. A product whose two inputs are both still being summed
is not a prescribed linear map.

If the coefficient maps are $B_j$ and the coordinate matrices $A_j$
satisfy
\[
 T_jB_{j-1}=B_jA_j,
\]
then
\[
 T_s\cdots T_1B_0=B_sA_s\cdots A_1.
\]
This follows by successive substitution. In particular a vector
in $\ker B_0$ remains physically zero after the contraction.
All phases and derivative coordinates are retained by the matrix
product. One may not discard a vector merely because the present
closing row kills it: a subsequent map can move that vector out of
the closing kernel.

No unitarity is assumed in this contraction. When estimates are
needed afterward, the incoming and outgoing Gram matrices must be
the transported pairings of Proposition~\ref{prop:coefficient-gram}.
For instance an isometry in these coordinates satisfies
$A_j^*G_jA_j=G_{j-1}$. An arbitrary multiplication by a fixed
tangle need not have this property.

\section{Fourier orthogonality on the shifted grid}

Fix $N\ge2$ and a real shift $\alpha$. Put
\[
 x_k=(k+\alpha)/N,\qquad
 w_k=e^{2\pi\ii x_k},\qquad 0\le k<N.
\]
The normalized Fourier matrix is
\[
 \mathcal U_{ak}=N^{-1/2}w_k^{-a},
                   \qquad0\le a,k<N.
\]
For these complete sets of indices,
\[
 \sum_{k=0}^{N-1}w_k^{a-b}
   =\begin{cases}N,&a=b,\\0,&a\ne b.\end{cases}
\]
Indeed, when $a\ne b$ the ratio
$e^{2\pi\ii(a-b)/N}$ is a nontrivial $N$th root of unity, so the
finite geometric sum vanishes. The factor from $\alpha$ is
constant in $k$ and does not change this conclusion. It follows
that $\mathcal U\mathcal U^*=I$ and $\mathcal U$ is unitary.

\begin{proposition}[Parseval in the coefficient normalization]
\label{prop:shifted-parseval}
For $f_0,\ldots,f_{N-1}\in\CC$, define
$H_a=N^{-1}\sum_k f_kw_k^{-a}$. Then
\[
 f_k=\sum_{a=0}^{N-1}H_aw_k^a,\qquad
 \sum_a|H_a|^2=\frac1N\sum_k|f_k|^2.
\]
\end{proposition}
\begin{proof}
Expand the first sum and use character orthogonality. For the
second, expand the squared moduli as
$N^{-2}\sum_{k,l}f_k\overline{f_l}\sum_aw_l^aw_k^{-a}$.
Only $k=l$ survives, with factor $N$.
\end{proof}

The filling argument uses $\alpha=1/2$, so $w_k^N=-1$. It also
retains derivatives at double spectral roots. The following estimate
isolates the role of orthogonality in that calculation.

\begin{proposition}[A bound for the full Hermite coefficient vector]
\label{prop:hermite-coefficient-mass}
Let $n\in\ZZ$, let $|h_k|=1$, and put
\[
 d_k=n\bigl(k-(N-1)/2\bigr)h_k,\qquad
 H_a=\frac1N\sum_kh_kw_k^{-a},\qquad
 D_a=\frac1N\sum_kd_kw_k^{-a}.
\]
Define $a_0=H_0$ and
$a_a=((N-a)H_a+D_a)/N$ for $1\le a<N$.
Set $a_N=a_{N+1}=0$ and
$\beta_j=(-1)^j(a_j-a_{j+2})$ for $0\le j<N$. Then
\[
 \sum_{j=0}^{N-1}|\beta_j|
       \le2\sqrt N\left(1+\frac{|n|}{\sqrt{12}}\right).
\]
\end{proposition}
\begin{proof}
Parseval gives $\|H\|_2^2=1$ and
\[
 \|D\|_2^2
 =\frac{n^2}{N}\sum_{k=0}^{N-1}\bigl(k-(N-1)/2\bigr)^2
 =\frac{n^2(N^2-1)}{12}.
\]
The last equality follows by expanding the square and using
$\sum k=N(N-1)/2$ and
$\sum k^2=N(N-1)(2N-1)/6$.
The $H$ part of $a$ is a diagonal contraction. The $D$ part is
divided by $N$ and its zeroth coordinate is omitted. Hence
\[
 \|a\|_2\le1+\frac{\|D\|_2}{N}
             \le1+\frac{|n|}{\sqrt{12}}.
\]
The truncated shift $a_j\mapsto a_{j+2}$ has norm at most one,
so $\|\beta\|_2\le2\|a\|_2$. Cauchy--Schwarz gives
$\|\beta\|_1\le\sqrt N\|\beta\|_2$.
\end{proof}

In the physical application $h_k=(-1)^{nk}q^{nk(k+1)}$ with
$q=e^{\pi\ii/N}$. Chapter~\ref{ch:conway-current} proves that these
values, the displayed normal data and the coefficient reconstruction
give the full crossing power, including negative and odd powers.
Proposition~\ref{prop:hermite-coefficient-mass} estimates that exact
coefficient vector. It is not a replacement of a double root by
its value alone.

\section{Eliminating stationary variables}

Let $F(y,t)$ be holomorphic near $(y_0,0)\in\CC^k\times\CC^r$.
Assume $F_y(y_0,0)=0$ and $A=F_{yy}(y_0,0)$ is invertible.
The holomorphic implicit-function theorem supplies a unique local
function $y_*(t)$ with $y_*(0)=y_0$ and $F_y(y_*(t),t)=0$.
Set $\Phi(t)=F(y_*(t),t)$.

\begin{proposition}[The stationary Hessian]
\label{prop:stationary-schur}
On this stationary family,
\[
 Dy_*=-F_{yy}^{-1}F_{yt},\qquad
 \Phi_t=F_t,\qquad
 \Phi_{tt}=F_{tt}-F_{ty}F_{yy}^{-1}F_{yt}.
\]
Every block is evaluated at $(y_*(t),t)$. If
$F(y,t)=F(y,-t)$ for all nearby $y,t$, then $F_{yt}(y,0)=0$.
\end{proposition}
\begin{proof}
Differentiate $F_y(y_*(t),t)=0$ to obtain
$F_{yy}Dy_*+F_{yt}=0$. The first formula follows by inversion.
The chain rule gives
$\Phi_t=F_t+F_yDy_*=F_t$. Differentiating it once more and
substituting the first formula gives the Schur complement. Evenness
gives $F_t(y,0)=0$ identically in $y$, hence $F_{yt}(y,0)=0$.
\end{proof}

The derivatives do not determine the additive constant of $\Phi$.
Its complete value must be kept when comparing it with a volume.
Likewise stationary elimination does not determine whether the
original finite sum contains the resulting Gaussian term. That
requires the finite-lattice contour argument and the nonzero density
in Chapter~\ref{ch:conway-current}, followed by the estimate of the
whole filling sum.

Grounding a graph Laplacian is a different operation. If
$L=B^{\mathsf T}WB$ is a real weighted graph Laplacian with connected
positive support, fixing one coordinate to zero removes its constant
kernel. The reduced principal matrix is positive definite. Eliminating
other free coordinates after that step produces a Schur complement.
The shared-weight variational arguments later in this reference use
both ideas in these respective roles.
